\documentclass[11pt, a4paper]{article}

\usepackage[utf8]{inputenc}
\usepackage{geometry}
\geometry{margin=1in}
\usepackage{hyperref}
\usepackage{graphicx}
\usepackage{tabularx}
\usepackage{booktabs}
\usepackage{enumitem}
\usepackage{titlesec}
\usepackage{xcolor}
\usepackage{listings}

% Color definitions
\definecolor{primary}{RGB}{0, 102, 204}
\definecolor{secondary}{RGB}{102, 102, 102}

% Title formats
\titleformat{\section}{\Large\bfseries\color{primary}}{}{0em}{}[\titlerule]
\titleformat{\subsection}{\large\bfseries\color{secondary}}{}{0em}{}
\titleformat{\subsubsection}{\normalsize\bfseries}{}{0em}{}

\title{
    \vspace{-2cm}
    \textbf{\Huge Comprehensive Testing Infrastructure Report}\\
    \vspace{0.5cm}
    \large An Expert-Guided Multimodal AI Ecosystem
}
\author{Development Team}
\date{\today}

\begin{document}

\maketitle

\begin{abstract}
This document serves as the complete technical blueprint for the automated testing infrastructure deployed in the Medical AI Ecosystem Backend. It thoroughly examines the configurations, unit validations, system-wide integration checks, Role-Based Access Control (RBAC) matrices, and performance load thresholds implemented to ensure robust software delivery.
\end{abstract}

\tableofcontents
\newpage

\section{Testing Architecture \& Configuration}
The testing suite relies on \texttt{pytest}, \texttt{pytest-django}, and \texttt{Locust} as its core frameworks.

\subsection{Configuration Files}
\begin{itemize}
    \item \textbf{\texttt{pytest.ini}:} Defines the Django settings module (\texttt{medical\_ai\_backend.settings}), specifies test discovery paths, and registers custom markers such as \texttt{@pytest.mark.stress} for selective test execution.
    \item \textbf{\texttt{conftest.py}:} A centralized hub providing modular, reusable fixtures for the entire test suite. Fixtures include:
    \begin{itemize}
        \item Role-specific user factories (\texttt{doctor\_user}, \texttt{patient\_user}, \texttt{admin\_user}, \texttt{researcher\_user}).
        \item Authenticated API client factories dynamically generating JWT tokens.
        \item Pre-populated database states such as \texttt{sample\_case\_with\_segmentation} and \texttt{sample\_report}, ensuring tests do not repeatedly duplicate database setup boilerplate.
    \end{itemize}
\end{itemize}

\section{Unit Testing}
Unit tests isolate and validate discrete components, primarily focusing on Django models and serializers.

\subsection{Database Models (\texttt{test\_models.py})}
\begin{itemize}
    \item \textbf{\texttt{Case} Model:} Validates UUID primary key generation, default status assignments (\textit{pending}), and relationship mapping to the \texttt{User} table.
    \item \textbf{\texttt{MRIImage} \& \texttt{SegmentationResult} Models:} Validates cascading deletion. When a \texttt{Case} is destroyed, tests verify that all linked \texttt{MRIImage} files and generated \texttt{SegmentationResult} masks are completely purged.
    \item \textbf{\texttt{Report} Model:} Verifies the state machine transitions (e.g., from \textit{draft} to \textit{finalized}), and confirms that \texttt{edit\_count} and \texttt{reviewed\_at} timestamps update appropriately.
\end{itemize}

\subsection{Serializers (\texttt{test\_serializers.py})}
\begin{itemize}
    \item \textbf{Data Transformation:} Strictly enforces \textit{camelCase} (frontend) to \textit{snake\_case} (backend) mapping across all serializers.
    \item \textbf{Validation Logic:} Confirms that \texttt{RegisterSerializer} correctly intercepts duplicate emails, enforces minimum password lengths, and rejects password mismatches.
    \item \textbf{Role Escalation Prevention:} Confirms that \texttt{UpdateProfileSerializer} intercepts and blocks malicious payload attempts (e.g., a patient injecting \texttt{role='admin'} during a profile update).
    \item \textbf{Read-Only Fields:} Ensures protected attributes (like \texttt{created\_at} and \texttt{case\_id}) cannot be altered by client payload requests.
\end{itemize}

\subsection{System Health (\texttt{test\_health.py})}
Validates that the \texttt{/api/health/} endpoint correctly returns HTTP 200 statuses and dictionary payloads indicating system uptime, ensuring compatibility with container orchestration liveness probes.

\section{Integration Testing}
Integration tests enforce the interactions between the database, serializers, and HTTP views.

\subsection{Authentication \& Lifecycles (\texttt{test\_auth.py})}
\begin{itemize}
    \item \textbf{Registration \& Login:} Verifies token issuance (access/refresh JWT) upon successful login and explicitly tests failures for non-existent users or incorrect passwords.
    \item \textbf{Token Management:} Validates secure logout mechanics (token blacklisting).
    \item \textbf{Account Deletion:} Tests the newly integrated cascading account deletion endpoint, ensuring an associated user's cases, reports, and data are entirely wiped upon secure password verification.
    \item \textbf{Admin User Management:} Validates that administrative staff can list, retrieve, update, and deactivate users, and verifies that deactivated users are subsequently blocked from logging in.
\end{itemize}

\subsection{Multimodal Inference Pipelines (\texttt{test\_pipeline.py})}
\begin{itemize}
    \item \textbf{Case Management (\texttt{TestCasePipeline}):} Tests the full CRUD capabilities on cases, including file uploads via \texttt{SimpleUploadedFile}.
    \item \textbf{Inference Execution (\texttt{TestInferencePipeline}):} Triggers the \texttt{/api/inference/predict/} endpoint. \textit{Crucial implementation detail:} \texttt{InferenceEngine} is rigorously mocked. The tests use \texttt{os.makedirs} to simulate physical case directories in the temporary test database, allowing path resolution checks to succeed without touching physical hardware.
    \item \textbf{Report Generation (\texttt{TestReportPipeline}):} Mocks the LLM utility functions (\texttt{generate\_json\_descriptor} and \texttt{generate\_report\_from\_descriptor}) to return predefined text, then tests the API's ability to edit, track edit counts, and export PDF artifacts.
\end{itemize}

\subsection{Strict RBAC Matrix (\texttt{test\_rbac\_matrix.py})}
A systematic matrix that executes endpoints across five unique authorization states to guarantee data silos:
\begin{itemize}
    \item \textbf{Patients:} Verified that they cannot access user lists, cannot delete cases, and cannot view reports unlinked to their account.
    \item \textbf{Doctors:} Verified they can list and manage their own cases but are blocked from accessing administrative user lists.
    \item \textbf{Researchers:} Confirmed they can create and process cases but cannot view a doctor's patient cases.
    \item \textbf{Admins:} Confirmed unrestricted Read/Write access across the ecosystem.
\end{itemize}

\section{System (End-to-End) Workflows}
End-to-End tests mimic complex real-world user journeys (\texttt{test\_e2e\_workflow.py}).

\begin{itemize}
    \item \textbf{Full Doctor Workflow:} Simulates a doctor registering, creating a case, uploading exactly four NIfTI modalities (\texttt{t1, t1ce, t2, flair}), running the inference engine, invoking the LLM for report generation, finalizing edits, and requesting a PDF export—all in one uninterrupted sequence.
    \item \textbf{Patient Read-Only Workflow:} Confirms a patient can register, locate their assigned case via their email linkage, and view (but explicitly not edit) the finalized diagnostic report.
    \item \textbf{Account Lifecycle Workflow:} Registers a user, updates their profile data, changes their password, and finally deletes their account, confirming \texttt{HTTP 401 Unauthorized} on subsequent login attempts.
\end{itemize}

\section{Stress \& Load Testing Architecture}
To guarantee system stability in production environments, custom stress testing scripts evaluate the ecosystem under extreme concurrency.

\subsection{Concurrency Verification (\texttt{test\_stress.py})}
These tests utilize Python's threading library and \texttt{TransactionTestCase} to prevent SQLite transaction locking.
\begin{itemize}
    \item \textbf{Concurrent Action Tests:} \texttt{TestConcurrentRegistration}, \texttt{TestConcurrentCaseCreation}, and \texttt{TestConcurrentCaseListing} launch 10-20 threads simultaneously to hammer endpoints. Assertions verify that zero race conditions occur and that precisely 10 cases/users are spawned.
    \item \textbf{API Response Boundaries:} \texttt{TestAPIResponseTime} enforces strict SLAs, asserting that logins execute under 2.0s, and list fetches under 1.0s.
\end{itemize}

\subsection{Unmocked Physical Inference Stress (\texttt{test\_inference\_stress.py})}
This specialized script targets the physical PyTorch MoME+ model and bypasses all testing mocks.
\begin{itemize}
    \item \textbf{Fake Data Mode:} Pushes 40 minimal 348-byte \texttt{nibabel} headers across 10 threads to test the Django server's preprocessing and error-handling capabilities without deadlocking.
    \item \textbf{Real Data Mode (\texttt{--real-data}):} Incorporates a dynamic \texttt{get\_nifti\_bytes} loader. When passed, it seamlessly loads $40$ authentic $4.5$MB NIfTI files from the \texttt{BraTSGLI} dataset, mapping formats like \texttt{t1c} to \texttt{t1ce}. This mode serves to explicitly benchmark physical GPU memory (VRAM) consumption, CUDA bottlenecks, and actual segmentation completion times under heavy concurrent load.
\end{itemize}

\subsection{API Load Testing (\texttt{Locust})}
The \texttt{locustfile.py} script sustains high-volume HTTP traffic against a live backend server. 
\begin{itemize}
    \item \textbf{Doctor Nodes (Weight: 3):} Continually perform heavy Read/Write operations (fetching profiles, generating new cases, and submitting MRI file uploads).
    \item \textbf{Patient Nodes (Weight: 1):} Simulate lighter Read operations (fetching assigned reports and viewing dashboard stats).
    \item \textbf{Benchmarks:} Execution with 50 concurrent users (spawning at 10/sec) resulted in an aggregated 516 requests with a $0\%$ failure rate. Median API response times registered under $10$ms natively.
\end{itemize}

\section{Conclusion}
The Medical AI Ecosystem Backend testing infrastructure provides absolute confidence in the platform's reliability. From exhaustive data serialization and strict RBAC enforcement to physical multi-threaded GPU load benchmarks, the suite effectively guards against regressions and guarantees enterprise-grade stability.

\end{document}
